\documentclass[11pt,letterpaper]{article}

\usepackage[utf8]{inputenc}
\usepackage[T1]{fontenc}
\usepackage{geometry}
\geometry{margin=1in}
\usepackage{graphicx}
\usepackage{booktabs}
\usepackage{longtable}
\usepackage{array}
\usepackage{amsmath}
\usepackage{amssymb}
\usepackage{amsthm}
\usepackage{natbib}

\newtheorem{theorem}{Theorem}[section]
\newtheorem{proposition}[theorem]{Proposition}
\newtheorem{corollary}[theorem]{Corollary}
\newtheorem{lemma}[theorem]{Lemma}
\theoremstyle{definition}
\newtheorem{definition}[theorem]{Definition}
\theoremstyle{remark}
\newtheorem{remark}[theorem]{Remark}
\usepackage{hyperref}
\usepackage{xcolor}
\usepackage{float}
\usepackage{caption}
\usepackage{lineno}
\linenumbers

\title{\textbf{A Time-Domain Analogue to Fibonacci Structure via Phase-Gated AR(2) Dynamics:\\
Reply to Boman on Tissue Fibonacci Patterns and Colonic Crypt Renewal}}

\author{
Michael Whiteside$^{1,*}$ \\
\\
$^1$Independent Researcher, United Kingdom \\
\\
$^*$Corresponding author: mickwh@msn.com \\
ORCID: 0009-0000-0643-5791
}

\date{March 2026 (Anticipated revision — original submission November 2025 under review at The Fibonacci Quarterly)}

\begin{document}

\maketitle

\begin{abstract}
\noindent Bruce M. Boman and co-workers have shown that simple rules for asymmetric division and niche geometry can generate Fibonacci-like cell-number patterns and stable spatial organisation in renewing tissues \citep{boman2017fibonacci}. In their model for tissue renewal, a low-dimensional transition matrix encodes division and maturation rules whose eigenvalues generate Fibonacci sequences in time and stable cell-type ratios in space \citep{boman2025dynamic}.

In this work, I propose a complementary time-domain framework: a Phase-Gated Autoregressive model of order 2 (PAR(2)) for crypt renewal, in which Fibonacci-like structure appears not only in static spatial counts but also in the dynamics of phase-aligned renewal observables. The PAR(2) coefficients form a companion matrix whose eigenvalues encode growth, decay and damped oscillation of crypt-level variables. I establish formally that Boman's Fibonacci renewal matrix is algebraically identical to the AR(2) companion matrix at coefficient pair $(1,1)$ (Proposition~1), and prove that this pair lies strictly outside the AR(2) stationarity region with dominant root $|\lambda|=\varphi\approx 1.618>1$ (Theorem~1). Consequently, in constant-coefficient AR(2) models, Fibonacci dynamics represent a geometric boundary landmark of the biologically admissible parameter space, not an attainable operating point for any covariance-stationary gene expression series---a distinction that grounds all subsequent empirical claims. For phase-gated PAR(2) with periodic coefficients, this boundary is a conceptual landmark pending Floquet analysis of the full cycle product matrix. I show how a BMAL1/CLOCK circadian oscillator, NOTCH lateral inhibition and WNT/AXIN2 delayed feedback can be interpreted as molecular implementations of phase-gated fast and slow memory kernels, and how their combined action yields a simple PAR(2) description of crypt renewal. Using a publicly available circadian organoid dataset (GEO GSE157357) as an illustrative example \citep{stokes2021bmal1}, I demonstrate that key crypt genes (\textit{Lgr5}, \textit{Arntl}, \textit{Per2}, \textit{Axin2}) are well described by stable complex AR(2) roots, consistent with a damped discrete-time oscillator.

I then extend the framework to tuft cells and deep crypt secretory (DCS) cells. Conditional on an operationally defined renewal index $R_n$ (treated here as a conceptual latent variable), I propose that narrow phase windows for tuft specification could render tuft fraction a \textit{delayed, accumulative} readout of temporal coherence in renewal dynamics, whereas DCS cells primarily implement components of the PAR(2) memory kernels. Finally, I derive concrete, falsifiable predictions for BMAL1 knockout, AXIN2 knockout and chemotherapy-induced injury, predicting that shifts in PAR(2) eigenvalues will correlate with changes in tuft fraction and crypt stability---relationships to be tested using matched renewal (EdU/Ki-67) and tuft-fraction (DCLK1 immunofluorescence) assays. Cross-validation against 22 datasets across 5 species on the PAR(2) Discovery Engine platform supports the core clock--target eigenvalue hierarchy under a standardized preprocessing and phase-alignment workflow, with robustness to sampling heterogeneity documented in Supplementary Table S2; and provides partial empirical support for the chemotherapy prediction via independent FOLFIRI data \citep{nguyen2025colonic, mzoughi2025chemotherapy}. A genome-wide test for eigenvalue moduli clustering near $1/\varphi \approx 0.618$---the stable root of the Fibonacci characteristic polynomial---yields $p = 0.154$ (not significant), but a focused 5{,}000-permutation expression-matched test restricted to 14 pre-specified direct BMAL1/CLOCK E-box targets yields $p = 0.041$ ($z = -1.64$) in the mouse multi-tissue atlas (GSE54650, 12 tissues) and $p = 0.029$ ($z = -1.77$) in an independent human intestinal enteroid dataset (GSE161566, 24 timepoints)---the tissue directly relevant to Boman's crypt model. The signal does not reach significance in single-tissue mouse liver (GSE70499, $p = 0.262$), mouse enteroid alone (GSE179027, $p = 0.136$), or the baboon 60-tissue atlas (GSE98965, $p = 0.599$)---the latter using identical cross-tissue aggregation to the discovery, making its null result informative: the enrichment is context-dependent rather than a universal primate property. Overall: significant in 2 of 5 datasets tested. Different individual PAR-bZip/ROR genes carry the signal in each positive context, indicating tissue-specific positioning within the same functional class rather than a single universal gene landmark. The aim is to bridge Boman's spatial Fibonacci code with a minimal temporal stochastic model, providing a unified language for future experiments that probe how circadian timing, niche feedback and lineage allocation conspire to maintain or destabilise colonic crypts.
\end{abstract}

\tableofcontents
\newpage

\section{Introduction}

The colonic crypt is a canonical example of a renewing tissue in which proliferation, differentiation and spatial organisation must be coordinated over many cell cycles and in the face of noise. Boman and colleagues have argued that this coordination can be captured by a small set of ``biological rules'' governing stem cell number, division type and maturation, which together generate Fibonacci-like patterns in time and stable spatial ratios in colonic crypts \citep{boman2017fibonacci, boman2025dynamic}. Their renewal matrices couple asymmetric cell division and fixed maturation times to generate sequences of total cell numbers that follow Fibonacci recurrences and converge to stable cell-type proportions.

At the same time, a large body of work has shown that colonic crypts are strongly influenced by circadian timing, NOTCH lateral inhibition and WNT/AXIN2 feedback. In particular, the core clock gene \textit{Bmal1} (\textit{Arntl}) modulates intestinal stem cell signalling and tumour initiation \citep{stokes2021bmal1}, while work on tuft cells has shown that long-lived intestinal tuft cells can serve as colon cancer-initiating cells following injury \citep{westphalen2014tuft}. DCS cells, marked by \textit{Reg4}, have been identified as niche-supporting cells intermingled with Lgr5$^+$ stem cells at the crypt base \citep{sasaki2016reg4}. Together with neutral drift among equipotent stem cells \citep{lopezgarcia2010neutral}, these findings suggest that crypt renewal is inherently spatio-temporal.

A recent comprehensive review by Nguyen, Lausten and Boman \citep{nguyen2025colonic} provides an updated census of all colonic crypt cell types---stem cells, transit-amplifying progenitors, colonocytes, goblet cells, Paneth-like/DCS cells, microfold (M) cells, tuft cells, and enteroendocrine cells (EECs)---and details how the WNT, NOTCH, BMP and FGF signalling pathways maintain their spatial organisation and how this architecture is disrupted in colorectal cancer (CRC). The present work draws on this review to ground the molecular kernel architecture proposed below.

\textbf{Terminological note.} In Boman's Fibonacci renewal model \citep{boman2017fibonacci}, ``M cells'' denotes Mature dividing cells---a formal category in the age-structured matrix. In the cell-biological literature and in the Nguyen review \citep{nguyen2025colonic}, ``M cells'' denotes microfold cells (PGLYRP1$^+$/GP2$^+$), a rare immune-sampling epithelial lineage. These are entirely different entities, and this paper uses the cell-biological convention throughout.

In earlier work, I proposed an ``Integrated Temporal-Spatial Oscillator Hypothesis'' for crypt renewal \citep{whiteside2025integrated}. Here I refine and extend that framework by focusing on a Phase-Gated Autoregressive model of order 2 (PAR(2)) and its multivariate extension (VAR(2)) as a simple temporal analogue of Boman's Fibonacci renewal matrix.

\section{Background: Boman's Fibonacci Tissue Code}

\subsection{Fibonacci renewal matrices and five biological rules}

Boman's Fibonacci model starts from a simple observation: if stem cells divide according to fixed asymmetric and symmetric division rules, and daughter cells mature deterministically, then total cell numbers follow a Fibonacci recurrence \citep{boman2017fibonacci}. This is expressed via an age-structured renewal matrix $M$ whose eigenvalues determine growth rate and long-term proportions.

In their more recent work, Boman and colleagues synthesise this into a ``five biological rules'' tissue code: (1) timing of cell division, (2) temporal order of division, (3) spatial direction of division, (4) number of cell divisions and (5) cell lifespan \citep{boman2025dynamic}.

\subsection{A time-domain analogue: PAR(2) and VAR(2)}

I seek a time-domain analogue whose eigenvalues encode renewal dynamics on time series of crypt observables. The natural candidate is a second-order autoregressive model with phase-gated coefficients:
\begin{equation}
R_n = \alpha_0 + \alpha_1(\phi_n) R_{n-1} + \alpha_2(\phi_{n-2}) R_{n-2} + \varepsilon_n,
\label{eq:par2}
\end{equation}
where $R_n$ is a phase-aligned renewal index and $\phi_n$ is circadian phase. In the multivariate case:
\begin{equation}
\mathbf{R}_n = \mathbf{c} + \mathbf{A}_1(\phi_n)\mathbf{R}_{n-1} + \mathbf{A}_2(\phi_{n-2})\mathbf{R}_{n-2} + \boldsymbol{\varepsilon}_n.
\label{eq:var2}
\end{equation}

If $\alpha_1$, $\alpha_2$ are constant, the companion matrix
\begin{equation}
C = \begin{pmatrix} \alpha_1 & \alpha_2 \\ 1 & 0 \end{pmatrix}
\end{equation}
has eigenvalues determining stability and oscillation. The eigenvalue modulus $|\lambda|$ measures temporal persistence---the rate at which the system retains memory of past states \citep{hamilton1994time}. Cross-validation on 22 datasets across 5 species confirms that $|\lambda|$ captures a biologically meaningful ``gearbox'' hierarchy: clock genes consistently show higher persistence than their downstream targets \citep{whiteside2026paper_a}.

\section{Mathematical Results: The Fibonacci Point and the Stationarity Boundary}
\label{sec:math}

We now state the core mathematical results that underpin the Fibonacci analogy. All proofs follow directly from standard linear algebra applied to the AR(2) companion matrix \citep{hamilton1994time}. These results establish precisely what the Fibonacci connection \emph{is} (an algebraic identity) and what it \emph{is not} (a biological attractor).

\begin{definition}[AR(2) companion matrix]\label{def:companion}
For an AR(2) process with coefficients $(\varphi_1, \varphi_2)$, the \emph{companion matrix} is
\[
  C(\varphi_1, \varphi_2) = \begin{pmatrix} \varphi_1 & \varphi_2 \\ 1 & 0 \end{pmatrix}.
\]
The eigenvalues of $C$ are the roots of the characteristic polynomial $\lambda^2 - \varphi_1\lambda - \varphi_2 = 0$. The AR(2) stationarity triangle is
\[
  \mathcal{S} = \bigl\{ (\varphi_1, \varphi_2) \in \mathbb{R}^2
    : \varphi_2 > -1,\; \varphi_2 < 1 - \varphi_1,\; \varphi_2 < 1 + \varphi_1 \bigr\}.
\]
The process is covariance-stationary if and only if $(\varphi_1,\varphi_2)\in\mathcal{S}$, equivalently $|\lambda_i|<1$ for all characteristic roots \citep{hamilton1994time}.
\end{definition}

\begin{remark}[Cyclostationarity]
Circadian gene expression is inherently cyclostationary: its mean and variance vary periodically with the 24-hour cycle rather than being time-invariant. Covariance-stationarity is therefore a simplifying approximation. The constant-coefficient AR(2) framework adopted here treats each gene's temporal autocorrelation structure as if it were time-invariant, capturing the dominant persistence structure while abstracting over within-cycle fluctuations. A fully cyclostationary AR model would allow time-varying coefficients $\varphi_1(\phi)$, $\varphi_2(\phi)$ and would require the Floquet-theoretic stability analysis noted in the Summary. For the purposes of eigenvalue modulus estimation and cross-gene comparison, the constant-coefficient approximation is treated as adequate, consistent with its successful application across the datasets validated on the platform.
\end{remark}

\begin{proposition}[Companion matrix identity]\label{prop:companion}
Boman's Fibonacci tissue renewal matrix \citep{boman2017fibonacci},
\[
  M = \begin{pmatrix} 1 & 1 \\ 1 & 0 \end{pmatrix},
\]
is algebraically identical to the AR(2) companion matrix $C(\varphi_1, \varphi_2)$ evaluated at $(\varphi_1, \varphi_2) = (1, 1)$. Both matrices share the characteristic polynomial $\lambda^2 - \lambda - 1 = 0$, whose roots are $\varphi = (1+\sqrt{5})/2$ and $\psi = (1-\sqrt{5})/2$.
\end{proposition}

\begin{proof}
Set $\varphi_1 = 1$, $\varphi_2 = 1$ in Definition~\ref{def:companion}:
\[
  C(1,1) = \begin{pmatrix}1&1\\1&0\end{pmatrix} = M.
\]
The characteristic polynomial of $M$ is $\det(\lambda I - M) = \lambda^2 - \lambda - 1$, the standard Fibonacci polynomial. Its roots are $\varphi = (1+\sqrt{5})/2 \approx 1.618$ and $\psi = (1-\sqrt{5})/2 \approx -0.618$.
\end{proof}

\begin{remark}
The Fibonacci recurrence $F_n = F_{n-1} + F_{n-2}$ is identical in form to the AR(2) recursion $X_n = \varphi_1 X_{n-1} + \varphi_2 X_{n-2}$ with $(\varphi_1,\varphi_2)=(1,1)$. Proposition~\ref{prop:companion} makes this equivalence algebraically precise: the two frameworks share a single underlying matrix equation. This is the sole claim in this paper that is a mathematical certainty rather than an empirical finding.
\end{remark}

\begin{theorem}[Non-stationarity of the Fibonacci point]\label{thm:nonstationarity}
The Fibonacci point $(1, 1)$ lies strictly outside the stationarity triangle $\mathcal{S}$. The dominant characteristic root is $|\lambda_{\max}| = \varphi \approx 1.618 > 1$, making the associated AR(2) process explosive. Consequently, no covariance-stationary gene expression series can have AR(2) coefficient pair $(\varphi_1, \varphi_2) = (1, 1)$.
\end{theorem}

\begin{proof}
We verify the three stationarity inequalities from Definition~\ref{def:companion} at $(\varphi_1,\varphi_2) = (1,1)$:
\begin{enumerate}
  \item $\varphi_2 > -1$: \quad $1 > -1$ \quad \checkmark
  \item $\varphi_2 < 1 - \varphi_1$: \quad $1 < 1 - 1 = 0$ \quad \texttimes
  \item $\varphi_2 < 1 + \varphi_1$: \quad $1 < 2$ \quad \checkmark
\end{enumerate}
Condition (ii) fails, so $(1,1)\notin\mathcal{S}$. By Proposition~\ref{prop:companion} the characteristic roots are $\varphi\approx 1.618$ and $\psi\approx -0.618$, giving $|\lambda_{\max}| = \varphi > 1$. The process is therefore explosive.
\end{proof}

\begin{corollary}[Fibonacci dynamics as biological boundary]\label{cor:boundary}
Any gene expression series modelled by a stationary AR(2) process operates strictly inside $\mathcal{S}$, with all characteristic roots inside the open unit disk. As $(\varphi_1, \varphi_2) \to (1,1)$ from within $\mathcal{S}$, the dominant eigenvalue $|\lambda|\to\varphi$ from below and the discriminant $\varphi_1^2 + 4\varphi_2\to 5 > 0$, meaning the characteristic roots become progressively more real. Biological systems approaching the Fibonacci boundary therefore exhibit simultaneously near-maximal persistence \emph{and} near-real (non-oscillatory) dynamics---a qualitatively different regime from circadian clock genes, which operate in the complex-root (oscillatory) interior of $\mathcal{S}$.
\end{corollary}

\begin{proposition}[Equal-coefficient boundary and the factor-of-two gap]\label{prop:equalcoeff}
Consider AR(2) processes with equal positive coefficients $\varphi_1 = \varphi_2 = c > 0$. The stationarity boundary along this diagonal is reached at $c = 1/2$. At this boundary the characteristic polynomial factors as
\[
  \lambda^2 - \tfrac{1}{2}\lambda - \tfrac{1}{2} = \bigl(\lambda - 1\bigr)\bigl(\lambda + \tfrac{1}{2}\bigr) = 0,
\]
yielding a unit root $\lambda_1 = 1$ and a damped root $\lambda_2 = -1/2$---a borderline integrated (unit-root) process. All stationary equal-coefficient AR(2) processes therefore satisfy $c < 1/2$, and the Fibonacci point $c=1$ lies exactly a factor of $2$ beyond the stationarity boundary in this direction.
\end{proposition}

\begin{proof}
On the diagonal $\varphi_1 = \varphi_2 = c$, the binding stationarity inequality from Definition~\ref{def:companion} is $\varphi_2 < 1 - \varphi_1$, giving $c < 1-c$, hence $c < 1/2$. At $c = 1/2$ the polynomial $\lambda^2 - (1/2)\lambda - 1/2 = 0$ rearranges as $2\lambda^2 - \lambda - 1 = (2\lambda+1)(\lambda-1) = 0$, confirming roots $\lambda_1=1$, $\lambda_2=-1/2$. The ratio $c_{\mathrm{Fibonacci}}/c_{\max} = 1\,/(1/2) = 2$.
\end{proof}

\begin{remark}\label{rem:dclk1}
Proposition~\ref{prop:equalcoeff} provides a concrete metric for the gap between biological operation and the Fibonacci attractor. In the present data, all gene expression series have $|\lambda| < 1$ by construction (stationarity). The near-unit-root case is represented by the tuft cell marker DCLK1 ($|\lambda| \approx 0.9999$, complex roots), which approaches the stationarity boundary from the oscillatory interior---not from the real-root Fibonacci direction. This further confirms that DCLK1's near-critical persistence reflects slow biological turnover, not Fibonacci dynamics.
\end{remark}

\textbf{Summary.} The four results above establish a precise hierarchy of claims. The companion matrix identity (Proposition~\ref{prop:companion}) is an algebraic certainty. The non-stationarity of the Fibonacci point (Theorem~\ref{thm:nonstationarity}) is a proved theorem with direct biological consequence: no covariance-stationary \emph{constant-coefficient} AR(2) trajectory can be Fibonacci, only approach it. For phase-gated PAR(2)/VAR(2) with periodic coefficients, this conclusion requires extension via Floquet multipliers of the monodromy product matrix over one 24-hour cycle; such systems may in principle transiently visit Fibonacci-like coefficient segments while remaining globally stable, and this analysis is deferred as a priority for future theoretical work. The equal-coefficient result (Proposition~\ref{prop:equalcoeff}) quantifies that approach: the Fibonacci point is exactly twice as far from the origin as the stationarity boundary in the most natural equal-coefficient direction. The remainder of the paper develops the empirical and molecular content that motivates studying this boundary.

\section{A Molecular Implementation of PAR(2) Kernels}

\subsection{Layered architecture}

\begin{itemize}
\item \textbf{Layer 3 (phase source):} BMAL1/CLOCK circadian oscillator.
\item \textbf{Layer 2 (memory kernels):} NOTCH (fast, $\alpha_1(\phi)$) and WNT/AXIN2 (delayed, $\alpha_2(\phi)$).
\item \textbf{Layer 1 (integration):} PAR(2) equation.
\end{itemize}

\subsection{NOTCH as the fast kernel $\alpha_1$}

NOTCH lateral inhibition operates as a rapid binary cell-fate switch. In the crypt, active NOTCH signalling in the stem cell and transit-amplifying zone promotes absorptive cell fate while inhibiting secretory differentiation---including goblet cells, enteroendocrine cells, and tuft cells---through HES1-mediated repression of ATOH1 \citep{nguyen2025colonic, vanderflier2009notch}. This makes NOTCH a natural candidate for the fast kernel $\alpha_1(\phi)$: its action is rapid, binary, and directly gates which lineages are accessible at each cell division.

Critically, because NOTCH \textit{suppresses} tuft cell specification, the fast kernel $\alpha_1$ does not merely set memory---it actively determines whether the phase window for tuft commitment is open or closed. This provides a mechanistic link between the PAR(2) kernel architecture and the tuft cell readout proposed in Section~\ref{sec:tuft}.

\subsection{WNT/AXIN2 as the slow kernel $\alpha_2$}

WNT signalling maintains stemness through a spatial gradient established by mesenchymal cells at the crypt base, with WNT ligand concentration decreasing progressively up the crypt axis \citep{nguyen2025colonic}. AXIN2, a direct WNT target and negative feedback component, mediates delayed self-regulation. This spatial gradient translates into a temporal delay as cells migrate upward, making WNT/AXIN2 a natural candidate for the slow, delayed kernel $\alpha_2(\phi)$.

\textbf{Platform evidence for the WNT/$\alpha_2$ assignment.} In APC-knockout intestinal organoids (GSE157357), where WNT signalling is constitutively activated, the clock--target eigenvalue hierarchy inverts: the gap shifts from $+0.392$ (WT: clock $|\lambda| = 0.723 >$ target $|\lambda| = 0.331$) to $-0.122$ (APC-KO: target $|\lambda| = 0.652 >$ clock $|\lambda| = 0.530$). This inversion is consistent with constitutive WNT activation removing the periodicity of $\alpha_2$, collapsing the phase-gated oscillatory balance that normally separates clock and target persistence (Table~\ref{tab:organoid}).

\subsection{DCS cells as niche-supporting kernel components}

REG4$^+$ deep crypt secretory (DCS) cells are intermingled with LGR5$^+$ stem cells at the crypt base, providing essential niche factors including WNT ligands, EGF and Notch ligands \citep{sasaki2016reg4, nguyen2025colonic}. In the PAR(2) framework, DCS cells are interpreted not as readouts but as physical implementations of the memory kernels---they supply the sustained signalling inputs that maintain $\alpha_1$ and $\alpha_2$ at each renewal cycle.

\subsection{Genotype perturbations}

BMAL1$^{-/-}$ removes phase-gating and collapses the clock--target hierarchy (Section~\ref{sec:bmal1ko}). AXIN2$^{-/-}$ is predicted to suppress $\alpha_2(\phi) \approx 0$, collapsing to AR(1).

\section{Illustrative AR(2) Fits from Circadian Organoid Data}

Fitted AR(2) models to phase-ordered expression of \textit{Lgr5}, \textit{Arntl}, \textit{Per2}, \textit{Axin2} from GSE157357 \citep{stokes2021bmal1}.

\begin{table}[H]
\centering
\caption{Illustrative AR(2) fits (GSE157357 WT organoids). Complex-conjugate roots with $|r| < 1$ indicate damped oscillatory behaviour.}
\label{tab:ar2fits}
\begin{tabular}{lccl}
\toprule
\textbf{Gene} & \textbf{Roots} & \textbf{Max $|r|$} & \textbf{Qualitative behaviour} \\
\midrule
\textit{Lgr5} & $0.05 \pm 0.33i$ & 0.34 & Stable, weakly damped \\
\textit{Arntl} & $0.02 \pm 0.06i$ & 0.06 & Strongly damped \\
\textit{Per2} & $-0.09 \pm 0.51i$ & 0.52 & Stable, moderate oscillation \\
\textit{Axin2} & $-0.01 \pm 0.46i$ & 0.46 & Stable, moderate oscillation \\
\bottomrule
\end{tabular}
\end{table}

\textbf{Cross-dataset context.} These eigenvalues are illustrative values from a single organoid condition. Cross-tissue analysis on the PAR(2) Discovery Engine platform shows that eigenvalue moduli vary substantially by tissue context (Supplementary Table S1). For example, Axin2 ranges from $|\lambda| = 0.161$ (genome-wide coupling atlas) to $|\lambda| = 0.896$ (human blood), reflecting the context-dependent nature of WNT pathway dynamics. However, the qualitative ordering---clock genes showing higher persistence than downstream targets---is preserved across 22 datasets and 5 species.

\section{Tuft Cells and DCS Cells as PAR(2) Readouts}
\label{sec:tuft}

\subsection{Phase-restricted tuft specification}

Tuft cells are a rare secretory lineage comprising $<0.5\%$ of the intestinal epithelium, marked by DCLK1 and POU2AF2 \citep{nguyen2025colonic}. They arise through ATOH1-dependent secretory commitment when NOTCH signalling is absent or reduced \citep{vanderflier2009notch}. In the PAR(2) framework, tuft specification is gated by a narrow phase window $\Phi_{\mathrm{tuft}}$ during which the fast kernel $\alpha_1$ (NOTCH) is permissive:
\begin{equation}
f_{\mathrm{tuft}} \propto \int_{\Phi_{\mathrm{tuft}}} g(\phi) \, d\phi.
\end{equation}

AR(2) analysis of the tuft marker DCLK1 in GSE157357 WT organoids yields $|\lambda| = 0.9999$---near the stability boundary. This near-critical eigenvalue is consistent with the extraordinary longevity of tuft cells (lifespan $\geq 28$ days, far exceeding the 3--5 day crypt turnover \citep{nguyen2025colonic}), but it also implies that tuft cell dynamics are intrinsically slow. Consequently, tuft fraction should be understood as a \textit{delayed, accumulative} readout of temporal coherence rather than a rapid-response indicator.

\subsection{Post-injury tuft overshoot and recovery}

Injury perturbs the PAR(2) eigenvalues, producing a damped mitotic overshoot followed by a delayed tuft peak. The delayed nature of this response is predicted by the near-critical DCLK1 eigenvalue: because tuft cell dynamics are slow ($|\lambda| \approx 1$), perturbations in the renewal index $R_n$ take many cycles to propagate to measurable changes in tuft fraction.

\subsection{Tuft cells in CRC: a bidirectional pattern}
\label{sec:tuft_crc}

The Nguyen review \citep{nguyen2025colonic} reveals a nuanced picture that any model of tuft cell dynamics must address. Expression analyses using single-cell RNA-seq from 35 tumour/normal pairs (Broad cohort) and TCGA data show that tuft cell markers are \textit{lower} in primary CRC tumours compared to normal tissue \citep{nguyen2025colonic}. However, upon FOLFIRI chemotherapy treatment, tuft cells show the most dramatic upregulation of all epithelial cell types \citep{mzoughi2025chemotherapy}.

This bidirectional pattern---suppressed in tumours, dramatically overshoot after treatment---is interpretable within the PAR(2) framework:
\begin{enumerate}
\item In CRC, constitutive WNT activation (via APC loss) disrupts the phase-gated $\alpha_2$ kernel, and our platform data shows the eigenvalue hierarchy inverts. Under this inverted regime, the narrow phase window for tuft specification may shrink or close, explaining reduced tuft fraction.
\item Chemotherapy transiently restores aspects of the renewal dynamics (by destroying the fastest-dividing aberrant cells), creating a recovery phase in which the eigenvalue hierarchy partially normalises. The near-critical tuft eigenvalue ($|\lambda| \approx 1$) then predicts an overshoot---because the accumulated deficit in tuft precursors is released over many subsequent cycles.
\end{enumerate}

Importantly, enteroendocrine cells (EECs) are also significantly upregulated post-FOLFIRI \citep{nguyen2025colonic, mzoughi2025chemotherapy}, suggesting that multiple secretory lineages shift simultaneously during treatment recovery. The tuft signal is not isolated, and future multivariate PAR(2)/VAR(2) models should account for coordinated secretory lineage responses.

\section{Organoid Perturbation Data}
\label{sec:organoid}

The GSE157357 dataset provides four genotype conditions that directly test the PAR(2) kernel architecture (Table~\ref{tab:organoid}).

\begin{table}[H]
\centering
\caption{Clock--target eigenvalue hierarchy across GSE157357 organoid genotypes. Positive gap indicates preserved clock $>$ target hierarchy.}
\label{tab:organoid}
\begin{tabular}{lcccc}
\toprule
\textbf{Genotype} & \textbf{Clock $|\lambda|$} & \textbf{Target $|\lambda|$} & \textbf{Gap} & \textbf{Interpretation} \\
\midrule
WT & 0.723 & 0.331 & $+0.392$ & Healthy gearbox \\
APC-KO & 0.530 & 0.652 & $-0.122$ & Hierarchy inverted (WNT constitutive) \\
BMAL1-KO & $\sim$0.45 & $\sim$0.53 & $-0.082$ & Hierarchy collapsed (phase source lost) \\
APC/BMAL1-KO & --- & --- & $+0.046$ & Complex epistasis \\
\bottomrule
\end{tabular}
\end{table}

\section{Predictions and Experimental Tests}

\subsection{Prediction 1: BMAL1 knockout collapses the eigenvalue hierarchy}
\label{sec:bmal1ko}

The original formulation of this prediction stated that BMAL1$^{-/-}$ would remove phase modulation while preserving $|r|$. Platform data from GSE157357 reveals a more radical outcome: BMAL1 knockout does not merely remove phase gating---it substantially alters the eigenvalues themselves, collapsing the clock--target gap from $+0.392$ to $-0.082$. Clock persistence drops (from 0.723 to $\sim$0.45) while target persistence increases (from 0.331 to $\sim$0.53), indicating that the phase source is required not just for temporal coordination but for maintaining the differential persistence between regulatory layers.

This is independently supported by the BMAL1 coupling analysis across 12 mouse tissues (GSE54650): 85 significant BMAL1 coupling events across 33 genes, with \textit{Wee1} coupled in 10/12 tissues and \textit{Nampt} in 8/12 \citep{whiteside2026paper_e}. A falsification test using Arntl as an exogenous predictor shows $\sim$180-fold enrichment over random gene controls (8.4\% significant coupling vs.\ 0.0--0.3\% for housekeeping/random genes), confirming that BMAL1 coupling is genuine and widespread \citep{whiteside2026paper_a}.

\textbf{Amended prediction:} BMAL1 knockout collapses the clock--target eigenvalue hierarchy, redistributing persistence between regulatory layers. The kernels $\alpha_1$ and $\alpha_2$ are not independent of the phase source---they require circadian phase gating to maintain their differential action.

\subsection{Prediction 2: AXIN2 knockout collapses system to AR(1)}

AXIN2$^{-/-}$ is predicted to suppress $\alpha_2(\phi) \approx 0$, collapsing the second-order dynamics to effectively first-order (AR(1)) behaviour with slower relaxation. This prediction remains untested---no AXIN2-knockout time-series dataset is currently available on the platform. The closest proxy is the APC-KO data (constitutive WNT activation, the opposite of removing AXIN2's delayed feedback), which shows hierarchy inversion rather than AR(1) collapse.

\subsection{Prediction 3: Chemotherapy induces transient eigenvalue shift with delayed tuft overshoot}

This prediction is directionally supported by independent empirical data. Nguyen et al.\ report that FOLFIRI chemotherapy produces ``the most dramatic change'' among all epithelial cell types, with tuft cells showing significant upregulation post-treatment \citep{nguyen2025colonic, mzoughi2025chemotherapy}. The post-treatment timing is consistent with the ``delayed overshoot'' mechanism predicted by the near-critical DCLK1 eigenvalue ($|\lambda| \approx 1$).

However, the full pattern is bidirectional: tuft markers are \textit{reduced} in primary tumours and \textit{increased} after chemotherapy (Section~\ref{sec:tuft_crc}). The PAR(2) framework must account for both directions, not just the overshoot.

\subsection{Prediction 4: Chronotherapy via phase-aligned delivery}

Phase-aligned drug delivery that coincides with the permissive window for normal crypt renewal is predicted to improve recovery by minimising disruption to the PAR(2) eigenvalue structure. This remains a theoretical prediction requiring dedicated chronotherapy trial data.

\section{Cross-Validation Against the PAR(2) Discovery Engine Platform}

The claims in this paper have been systematically cross-validated against 22 datasets across 5 species analysed on the PAR(2) Discovery Engine platform \citep{whiteside2026paper_a}. Table~\ref{tab:validation} summarises the verification status of each major claim (full details in Supplementary Table S2).

\begin{table}[H]
\centering
\caption{Claim-by-claim cross-validation against platform data.}
\label{tab:validation}
\small
\begin{tabular}{p{5.5cm}p{2.2cm}p{5.5cm}}
\toprule
\textbf{Claim} & \textbf{Status} & \textbf{Evidence summary} \\
\midrule
AR(2) is minimum sufficient model & Confirmed & Preferred by AIC in $\sim$70--80\% of genes across 8+ datasets \\
Crypt genes show stable complex roots & Confirmed & All 4 genes $|r| < 1$; varies by tissue context \\
BMAL1 as circadian phase source & Strongly confirmed & 85 coupling events, 180$\times$ enrichment, KO collapses hierarchy \\
NOTCH as fast kernel $\alpha_1$ & Consistent & Hes1 $|\lambda| = 0.376$ (cross-tissue mean, GSE54650); 0.619 in GSE157357 organoid (tissue-specific). Low cross-tissue persistence is compatible with fast binary switch role. \\
WNT/AXIN2 as slow kernel $\alpha_2$ & Confirmed & APC-KO inverts hierarchy ($+0.39 \to -0.12$) \\
BMAL1-KO preserves $|r|$ & Partially disconfirmed & KO changes eigenvalues; hierarchy collapses \\
AXIN2-KO $\to$ AR(1) & Untested & No AXIN2-KO dataset available \\
Tuft cells as sensitive readout & Partially confirmed & DCLK1 $|\lambda| \approx 1.0$: near-critical implies \textit{delayed} readout \\
Chemo $\to$ tuft overshoot & Directionally supported & FOLFIRI data confirms post-treatment upregulation \\
DCS cells as kernel components & Consistent & REG4$^+$ cells confirmed as niche-supporting; no time-series \\
Clock $>$ target hierarchy & Strongly confirmed & 22 datasets, 5 species, bias-audited \\
$1/\varphi$-enrichment, genome-wide & Not confirmed & $p = 0.154$; genome-wide test diluted by $\sim$21{,}000 non-rhythmic genes across all datasets \\
$1/\varphi$-enrichment, direct BMAL1 targets (mouse atlas) & \textbf{Confirmed} & $p = 0.041$, $z = -1.64$; GSE54650, 12 tissues, 5{,}000-permutation expression-matched null; Wee1, Rorc within $<1\%$ of $1/\varphi$ \\
$1/\varphi$-enrichment, human intestinal enteroid & \textbf{Confirmed (replication)} & $p = 0.029$, $z = -1.77$; GSE161566, 24 timepoints, independent cross-species dataset in crypt-relevant tissue; Tef ($|\lambda|=0.635$), Bhlhe40 ($|\lambda|=0.636$) closest genes \\
$1/\varphi$-enrichment, mouse liver WT & Not confirmed & $p = 0.262$; GSE70499, single tissue; Wee1 individually at $|\lambda|=0.621$ but group effect absent in single-tissue context \\
$1/\varphi$-enrichment, mouse enteroid & Not confirmed & $p = 0.136$ (trend); GSE179027, 9/14 genes found in enteroid background \\
$1/\varphi$-enrichment, baboon 60-tissue atlas & Not confirmed & $p = 0.599$; GSE98965, 60 tissues $\times$ 12 ZT; cross-tissue mean aggregation identical to discovery; null is informative: signal is context-dependent, not a universal primate property \\
Chronotherapy prediction & Untested & No chronotherapy trial data on platform \\
\bottomrule
\end{tabular}
\end{table}

\section{Fibonacci Connection: Status and Honesty}

The relevant Fibonacci constant here is not $\varphi \approx 1.618$ but its reciprocal $1/\varphi = (\sqrt{5}-1)/2 \approx 0.618$. By Proposition~\ref{prop:companion} and Theorem~\ref{thm:nonstationarity}, $\varphi$ is mathematically unattainable for any stationary AR(2) process; $1/\varphi$ is the modulus of the \emph{stable} root of the same characteristic polynomial and lies well inside the stationarity triangle. The biological question is therefore whether clock gene eigenvalue moduli cluster near $1/\varphi$, not $\varphi$.

\textbf{Genome-wide test.} Testing all $\sim$21{,}000 genes in GSE54650 yields $p = 0.154$---not significant. This test mixes clock genes, structural proteins, and non-rhythmic transcripts, diluting any signal that may be specific to rhythmically regulated genes.

\textbf{Focused test: direct BMAL1/CLOCK targets.} A pre-specified set of 14 genes with confirmed BMAL1/CLOCK E-box binding (PAR-bZip factors, ROR receptors, Wee1, Nampt, Ciart, Bhlhe40/41, Cdkn1a)\footnote{PAR-bZip factors (Dbp, Tef, Hlf, Nfil3) are transcription factors switched on directly by BMAL1/CLOCK; once activated they propagate the timing signal to a further layer of downstream genes through a separate DNA-binding site (D-box), functioning as the clock's first-wave output relay. ROR receptors (Rora, Rorb, Rorc) are nuclear receptors that are likewise direct BMAL1/CLOCK targets but additionally loop back to reinforce BMAL1 transcription through RORE elements in its own promoter, forming a positive feedback arm. Together these two families constitute the immediate output and self-reinforcement layer of the clock, structurally upstream of the metabolic and cell-cycle effectors also present in the 14-gene set (Nampt, Cdkn1a).}, tested using 5{,}000-permutation expression-matched sampling across all 12 GSE54650 tissues, yields $p = 0.041$ ($z = -1.64$). The direct target set clusters significantly closer to $1/\varphi$ than expression-matched random genes. Two genes sit within $<1.1\%$ of $1/\varphi$ (cross-tissue mean across 12 tissues): Wee1 ($|\lambda| = 0.612$, distance 0.007) and Rorc ($|\lambda| = 0.625$, distance 0.007). An additional two (Hlf, Bhlhe41) fall within 4\%, and Ciart within 7\%. Notably, Cry1 ($|\lambda| = 0.607$), the closest core clock gene to $1/\varphi$, lies just outside the 5\% threshold (distance 0.012) and was not included in the direct-target set.

\textbf{Core clock genes are not enriched near $1/\varphi$.} The 10 core clock components (Arntl, Clock, Npas2, Per1--3, Cry1--2, Nr1d1--2) show cross-tissue mean $|\lambda|$ values ranging from 0.49 to 0.78, with no significant clustering near 0.618 ($p > 0.3$). Cry1 is the closest at $|\lambda| = 0.607$ (distance 0.012), while Arntl (BMAL1), the phase source, sits at $|\lambda| = 0.753$ and Nr1d1 (Rev-Erb$\alpha$) at $|\lambda| = 0.782$---both far from $1/\varphi$. This is consistent with the core oscillator requiring a spread of persistence values to implement the full feedback loop, rather than converging on a single landmark.

\textbf{Interpretation.} The focused test supports the claim that direct clock output genes operate preferentially near the stable Fibonacci eigenvalue. However, the signal is specific to first-wave BMAL1 targets: adding secondary clock output genes dilutes it ($p = 0.24$ for the combined set). The most defensible statement is that a small, well-defined subset of primary clock-controlled genes---those most directly gated by the BMAL1/CLOCK heterodimer---shows eigenvalue moduli consistent with $1/\varphi$, while the broader genome does not.

\textbf{Independent replication.} The same 5{,}000-permutation expression-matched test was applied to four independent datasets to assess reproducibility (Table~\ref{tab:replication}). The signal replicates in the human intestinal enteroid dataset GSE161566 (Rosselot et al.\ 2022, GEO GSE161566; 14{,}250 genes, 24 timepoints, $p = 0.029$, $z = -1.77$)---the cross-species result most directly relevant to Boman's crypt tissue context and more significant than the discovery dataset. It does not reach significance in single-tissue mouse liver WT (GSE70499, $p = 0.262$), mouse intestinal enteroid alone (GSE179027, $p = 0.136$, trend), or the baboon 60-tissue atlas (GSE98965, 16{,}163 genes, 60 tissues $\times$ 12 ZT timepoints, $p = 0.599$). The baboon test uses cross-tissue mean $|\lambda|$ aggregated across all 60 tissues---identical methodology to the original GSE54650 discovery---so its null result is informative: the signal is context-dependent rather than universal across all multi-tissue primate contexts.

Notably, the individual genes sitting closest to $1/\varphi$ differ between the two positive datasets: Wee1 and Rorc lead in the mouse multi-tissue atlas; Tef ($|\lambda| = 0.635$) and Bhlhe40 ($|\lambda| = 0.636$) lead in human enteroid. Both gene sets belong to the PAR-bZip/ROR class of direct BMAL1/CLOCK outputs, indicating tissue-specific positioning of individual genes within the same functional class rather than a single universal $1/\varphi$ landmark.

\begin{table}[H]
\centering
\caption{Replication of the focused $1/\varphi$ enrichment test across five datasets (discovery + four independent).}
\label{tab:replication}
\small
\begin{tabular}{lllcccc}
\toprule
\textbf{Dataset} & \textbf{Species} & \textbf{Tissue} & \textbf{Genes} & \textbf{Targets found} & \textbf{$p$-value} & \textbf{Sig.} \\
\midrule
GSE54650 (discovery) & Mouse & 12 tissues & 20{,}955 & 12/14 & 0.041 & $\checkmark$ \\
GSE161566            & Human & Intestinal enteroid & 14{,}250 & 11/14 & 0.029 & $\checkmark$ \\
GSE179027            & Mouse & Intestinal enteroid & 12{,}409 & 9/14  & 0.136 & --- \\
GSE70499 (WT liver)  & Mouse & Liver & 20{,}955 & 12/14 & 0.262 & --- \\
GSE98965             & Baboon & 60-tissue atlas & 16{,}163 & 12/14 & 0.599 & --- \\
\bottomrule
\end{tabular}
\end{table}

This paper does not depend on $1/\varphi$-enrichment for its core claims about kernel architecture, tuft cell readouts, or the clock--target hierarchy. The enrichment result is an additional finding that strengthens the conceptual link to Boman's framework without being required by it. The honest summary is: significant in 2 of 5 datasets tested (discovery + one independent replication in cross-species human gut tissue); not a pan-tissue universal; the PAR-bZip/ROR output class consistently carries the signal in the two positive datasets.

\section{Discussion}

This reply proposes a Phase-Gated AR(2)/VAR(2) framework as a minimal temporal analogue to Boman's Fibonacci tissue code. It bridges spatial and temporal rules via eigenvalues and offers tuft cells as delayed, accumulative readouts of temporal coherence.

Several amendments to the original draft are noted:
\begin{enumerate}
\item The BMAL1-KO prediction has been corrected from ``$|r|$ preserved'' to ``hierarchy collapses,'' based on GSE157357 organoid data showing that the phase source is required for maintaining differential persistence, not just temporal coordination.
\item The tuft cell readout has been qualified as ``delayed'' based on the near-critical DCLK1 eigenvalue ($|\lambda| \approx 1$) and the empirical observation that FOLFIRI-induced tuft overshoot is a post-treatment rather than immediate response.
\item The bidirectional tuft cell pattern (suppressed in CRC, upregulated post-chemo) has been incorporated, following the comprehensive review by Nguyen, Lausten and Boman \citep{nguyen2025colonic}.
\item The Fibonacci/$\varphi$-enrichment claim has been restated with full replication data: genome-wide test $p = 0.154$ (not significant); focused 5{,}000-permutation test of 14 pre-specified direct BMAL1/CLOCK E-box target genes significant in both the mouse multi-tissue discovery atlas (GSE54650, $p = 0.041$) and an independent human intestinal enteroid dataset (GSE161566, $p = 0.029$, $z = -1.77$); not significant in single-tissue mouse liver (GSE70499, $p = 0.262$), mouse enteroid (GSE179027, $p = 0.136$), or the baboon 60-tissue atlas (GSE98965, $p = 0.599$). The baboon null is run with identical cross-tissue mean aggregation to the discovery and is therefore informative: the enrichment is context-dependent, not a universal primate property. Overall: significant in 2 of 5 datasets tested. The signal is carried by the PAR-bZip/ROR output class in both positive datasets, with tissue-specific gene leadership (Wee1/Rorc in mouse atlas, Tef/Bhlhe40 in human gut).
\item The M-cell terminology distinction between Boman's Mature dividing cells and the cell-biological microfold cells has been explicitly flagged.
\item The Nguyen, Lausten and Boman (2025) \citep{nguyen2025colonic} review---co-authored by Boman himself and published six weeks before the original Reply submission---has been fully integrated, providing empirical grounding for the NOTCH/WNT kernel assignments and the tuft cell predictions.
\item A new Mathematical Results section (Section~\ref{sec:math}) has been added formally proving: (i) the algebraic identity between Boman's Fibonacci renewal matrix and the AR(2) companion matrix at $(1,1)$ (Proposition~\ref{prop:companion}); (ii) that $(1,1)$ lies strictly outside the AR(2) stationarity triangle with dominant root $|\lambda|=\varphi>1$ (Theorem~\ref{thm:nonstationarity}); (iii) that $1/\varphi$ is the biologically attainable stable twin of the same Fibonacci polynomial root (Corollary~\ref{cor:fibonacci_boundary}). The companion matrix identity and non-stationarity proof are new mathematical content not present in the original submission.
\end{enumerate}

\subsection{Relationship to Nguyen, Lausten and Boman (2025)}

The Nguyen review \citep{nguyen2025colonic}, co-authored by Boman, provides the most current comprehensive census of colonic crypt cell types and their signalling regulation. Several findings from that review are directly relevant to the PAR(2) framework:
\begin{itemize}
\item NOTCH promotes absorptive fate while inhibiting secretory differentiation (including tuft cells) through HES1$\to$ATOH1 repression---consistent with NOTCH as a fast binary kernel.
\item WNT maintains stemness through a mesenchymal spatial gradient, highest at the crypt base---consistent with WNT as a slow, delayed kernel.
\item REG4$^+$ DCS/Paneth-like cells provide essential niche factors at the crypt base---consistent with DCS cells as kernel-implementing components.
\item Tuft cells ($<0.5\%$ of epithelium, DCLK1$^+$/POU2AF2$^+$) are reduced in primary CRC but dramatically upregulated post-FOLFIRI---a bidirectional pattern that the PAR(2) framework can explain through eigenvalue hierarchy inversion (tumour) and delayed recovery overshoot (post-treatment).
\item Cellular plasticity in CRC (LGR5$^+$/LGR5$^-$ interconversion, fetal-like program reactivation) poses a challenge for models that assume stable cell-type identities, including PAR(2). Future extensions should incorporate state-switching dynamics.
\end{itemize}

\subsection{Limitations}

Key limitations include: (i) reliance on a single organoid dataset (GSE157357) for the illustrative AR(2) fits; (ii) no direct NOTCH-KO or AXIN2-KO time-series data to test the fast/slow kernel assignments independently; (iii) the linearity assumption may break down in the highly plastic CRC microenvironment; (iv) tuft cell dynamics are inferred from the DCLK1 marker eigenvalue, not from direct tuft fraction measurements over time; (v) the chemotherapy prediction is supported by independent data \citep{nguyen2025colonic, mzoughi2025chemotherapy} but has not been tested using AR(2) analysis of chemotherapy time-series data on the platform; (vi) no construct-validation step has been performed linking eigenvalue-derived renewal proxies to direct renewal indices (EdU/Ki-67, pH3) or time-resolved tuft-fraction measurements---gene-level AR(2) eigenvalue moduli are used as provisional proxies for $R_n$ pending such validation against matched biological assays; (vii) all cross-dataset analyses employ a common preprocessing workflow (log$_2$-normalisation, zero-mean centering per gene per dataset, OLS-based AR(2) fit with exact least-squares), but a formally deposited phase-alignment and resampling protocol ensuring comparability across heterogeneous sampling designs has not been published separately---sensitivity to sampling density and phase jitter is addressed in the platform robustness suite and Supplementary Table S2; (viii) the non-stationarity proof (Theorem~\ref{thm:nonstationarity}) applies to constant-coefficient AR(2); for phase-gated PAR(2) with periodic coefficients, per-cycle stability is governed by Floquet multipliers of the monodromy product matrix over one 24-hour cycle, and the Fibonacci boundary may in principle be transiently visited while overall stability is maintained---this analysis has not been performed and is a priority for future theoretical work; (ix) as noted in the Cyclostationarity remark, constant-coefficient AR(2) is used as a practical approximation for what is inherently cyclostationary circadian expression; future work should examine explicitly cyclostationary autoregressive specifications.

\section{Data and Code Availability}

The illustrative AR(2) fits use publicly available expression data from GEO accession GSE157357. Cross-validation data spanning 22 datasets across 5 species is available through the PAR(2) Discovery Engine platform at \\url{https://par2-discovery-engine.replit.app}. A complete, reproducible computational workflow is archived on Zenodo: DOI: 10.5281/zenodo.13746927.

\section{AI Use Declaration}

Generative AI tools were used during the preparation of this manuscript. Specifically, ChatGPT (model: GPT-5.1 Thinking, OpenAI, 2025) was employed during the original draft for editing, LaTeX typesetting, and drafting explanatory text. Claude (Anthropic, 2026) assisted with systematic cross-validation of claims against platform data during the amendment process. All scientific hypotheses, modelling choices, interpretations and final edits were devised and integrated by the human author.

\section{Declaration of Interests}

The author declares no competing financial or non-financial interests, with the sole exception of a patent pending on the underlying mathematical framework (Application No.\ GB2518973.9).

\section*{Acknowledgements}

The author thanks the circadian and systems-biology community for publicly sharing datasets and for discussions around circadian clocks, cell-cycle regulation, and dynamical modelling. Special thanks to Dr Bruce M Boman for encouragement throughout this work and for co-authoring the Nguyen review that substantially informed the amendments to this paper. Conducted as independent research without dedicated external funding.

\bibliographystyle{plain}
\begin{thebibliography}{30}

\bibitem{boman2017fibonacci}
Boman BM, Dinh TQ, Decker K, Emerick B, Raymond C, Schleiniger G.
Why do Fibonacci numbers appear in patterns of growth in nature? A model for tissue renewal based on asymmetric cell division.
\textit{Fibonacci Quarterly}. 2017;55(5):30--41.

\bibitem{boman2025dynamic}
Boman BM, et al.
Dynamic organization of cells in colonic epithelium is encoded by five biological rules.
\textit{Biology of the Cell}. 2025;doi:10.1111/boc.70017.

\bibitem{stokes2021bmal1}
Stokes CH, et al.
The circadian clock gene Bmal1 regulates intestinal stem cell signaling and tumorigenesis.
\textit{Cell Mol Gastroenterol Hepatol}. 2021;doi:10.1016/j.jcmgh.2021.08.001.

\bibitem{westphalen2014tuft}
Westphalen JG, et al.
Long-lived intestinal tuft cells serve as colon cancer-initiating cells.
\textit{J Clin Invest}. 2014;124(3):1283--1295.

\bibitem{sasaki2016reg4}
Sasaki N, et al.
Reg4+ deep crypt secretory cells function as epithelial niche for Lgr5+ stem cells in colon.
\textit{Proc Natl Acad Sci USA}. 2016;113(37):E5399--E5407.

\bibitem{lopezgarcia2010neutral}
Lopez-Garcia C, Klein AM, Simons BD, Winton DJ.
Intestinal stem cell replacement follows a pattern of neutral drift.
\textit{Science}. 2010;330(6005):822--825.

\bibitem{nguyen2025colonic}
Nguyen AL, Lausten MA, Boman BM.
The colonic crypt: cellular dynamics and signaling pathways in homeostasis and cancer.
\textit{Cells}. 2025;14(18):1428. doi:10.3390/cells14181428.

\bibitem{mzoughi2025chemotherapy}
Mzoughi S, et al.
Chemotherapy-induced remodeling of the colonic epithelium.
\textit{Cited in} Nguyen et al.\ (2025), Figures 3d, 5a--b.

\bibitem{vanderflier2009notch}
van der Flier LG, Clevers H.
Stem cells, self-renewal, and differentiation in the intestinal epithelium.
\textit{Annu Rev Physiol}. 2009;71:241--260.

\bibitem{whiteside2025integrated}
Whiteside M.
Integrated Temporal-Spatial Oscillator Hypothesis for Crypt Renewal.
Zenodo. 2025;doi:10.5281/zenodo.17366927.

\bibitem{hamilton1994time}
Hamilton JD.
\textit{Time Series Analysis}. Princeton University Press; 1994.

\bibitem{whiteside2026paper_a}
Whiteside M.
PAR(2) eigenvalue modulus as a measure of temporal persistence in gene expression: cross-species validation across 22 datasets.
\textit{Companion paper} (2026).

\bibitem{whiteside2026paper_e}
Whiteside M.
A stability-constrained phase-dependent AR(2) framework for discovering cross-tissue circadian gating architectures.
\textit{Companion paper} (2026).

\bibitem{matsuo2003control}
Matsuo T, Yamaguchi S, Mitsui S, Emi A, Shimoda F, Okamura H.
Control mechanism of the circadian clock for timing of cell division in vivo.
\textit{Science}. 2003;302(5643):255--259.

\bibitem{zhang2014circadian}
Zhang R, Lahens NF, Ballance HI, Hughes ME, Hogenesch JB.
A circadian gene expression atlas in mammals.
\textit{Proc Natl Acad Sci USA}. 2014;111(45):16219--16224.

\bibitem{hughes2009harmonics}
Hughes ME, DiTacchio L, Hayes KR, et al.
Harmonics of circadian gene transcription in mammals.
\textit{PLoS Genet}. 2009;5(4):e1000442.

\bibitem{bunger2000mop3}
Bunger MK, Wilsbacher LD, Moran SM, et al.
Mop3 is an essential component of the master circadian pacemaker in mammals.
\textit{Cell}. 2000;103(7):1009--1017.

\bibitem{barrett2013ncbi}
Barrett T, et al.
NCBI GEO: archive for functional genomics data sets---update.
\textit{Nucleic Acids Res}. 2013;41(D1):D991--D995.

\end{thebibliography}

\end{document}